\DocumentMetadata{
  pdfversion=2.0,
  pdfstandard=ua-2,
  lang=en-US,
  tagging=on,
  tagging-setup={math/setup=mathml-SE,tabsorder=structure}
}
\documentclass[11pt,letterpaper]{article}
\usepackage[margin=0.68in,includefoot]{geometry}
\usepackage{newcomputermodern}
\usepackage{amsmath,amscd,mathtools}
\usepackage{tikz,adjustbox}
\providecommand{\square}{\mdlgwhtsquare}
\usepackage[hidelinks]{hyperref}
\hypersetup{
  pdftitle={Numerical Analysis PhD Exam, September 2009},
  pdfauthor={University of Florida Department of Mathematics},
  pdfsubject={Graduate mathematics examination},
  pdfdisplaydoctitle=true,
  bookmarksopen=true
}
\setcounter{secnumdepth}{0}
\setlength{\parindent}{0pt}
\setlength{\parskip}{0.28em}
\setlength{\emergencystretch}{3em}
\setlength{\leftmargini}{2.15em}
\setlength{\leftmarginii}{2.65em}
\setlength{\leftmarginiii}{3em}
\pagestyle{empty}
\raggedbottom
\allowdisplaybreaks
\NewTaggingSocketPlug{block/list/label}{examproblem}{%
  \tagstructbegin{tag=Lbl}\tagstructend%
  \tagstructbegin{tag=\LBody}%
  \tagstructbegin{tag=\ProblemHeadingTag,title-o={\ProblemHeadingText}}%
  \tagmcbegin{tag=\ProblemHeadingTag,actualtext={\ProblemHeadingText}}%
  #2%
  \tagmcend%
  \tagstructend%
}
\newenvironment{examproblems}[2]{%
  \def\ProblemHeadingTag{#2}%
  \AssignTaggingSocketPlug{block/list/label}{examproblem}%
  \begin{enumerate}%
  \setcounter{enumi}{#1}%
  \renewcommand{\labelenumi}{\textbf{\theenumi.}}%
  \setlength{\topsep}{0.35em}%
  \setlength{\partopsep}{0pt}%
  \setlength{\itemsep}{0.48em}%
  \setlength{\parsep}{0.18em}%
}{\end{enumerate}}
\newenvironment{examsubparts}{%
  \AssignTaggingSocketPlug{block/list/label}{default}%
  \begin{enumerate}%
  \setlength{\topsep}{0.18em}%
  \setlength{\partopsep}{0pt}%
  \setlength{\itemsep}{0.16em}%
  \setlength{\parsep}{0.1em}%
}{\end{enumerate}}
\newenvironment{examinstructions}{%
  \par\begingroup\itshape
}{%
  \par\endgroup\vspace{0.18em}
}
\makeatletter
\renewcommand\subsection{\@startsection{subsection}{2}{0pt}%
  {-1.25ex plus -.25ex minus -.1ex}%
  {0.55ex plus .1ex}%
  {\normalfont\large\bfseries}}
\renewcommand{\@maketitle}{%
  \newpage\null\vskip -1em%
  {\footnotesize\bfseries
    \href{https://www.ufl.edu/}{University of Florida}, \href{https://math.ufl.edu/}{Department of Mathematics}\par}%
  \vskip -0.5em%
  \tagstructbegin{tag=H1,title={Numerical Analysis PhD Exam, September 2009}}%
  \tagpdfparaOff%
  \tagmcbegin{tag=H1}%
    {\LARGE\bfseries\href{https://gma.math.ufl.edu/exams/numerical-analysis-phd/}{Numerical Analysis PhD Exam}, September 2009\par}%
  \tagmcend%
  \tagpdfparaOn%
  \tagstructend%
  \vskip 0.25em%
  \tagmcbegin{artifact}%
    \hrule height 1pt%
  \tagmcend%
  \vskip 1em%
}
\makeatother
\title{Numerical Analysis PhD Exam, September 2009}
\author{}
\date{}
\begin{document}
\maketitle
\thispagestyle{empty}
\begin{examinstructions}
Answer any 8 questions.
\end{examinstructions}
\begin{examproblems}{0}{H2}
\def\ProblemHeadingText{Problem 1.}
\item
This problem has two parts.
\begin{examsubparts}
\item[\textnormal{a.}]
Suppose~\(u\) and \(v\in\mathbb C^n\). Derive a formula for the~\(p\)-norm of \(uv^*\), \(p\ge 1\).
\item[\textnormal{b.}]
Let \(A=U\Sigma V^*\) be the singular value decomposition of~\(A\), where the diagonal elements of~\(\Sigma\) are in decreasing order:
\[
\sigma_1\geq\sigma_2\geq\sigma_3\geq\cdots.
\]
 Show that \(\|A\|_2=\sigma_1\) and
\[
\|A\|_F=\sqrt{\sigma_1^2+\sigma_2^2+\sigma_3^2+\cdots}.
\]
\end{examsubparts}
\def\ProblemHeadingText{Problem 2.}
\item
If \(A=QR\) is the \(QR\) factorization of~\(A\), show that \(r_{ii}\) is the distance from the~\(i\)-th column of~\(A\) to the space spanned by the first \(i-1\) columns of~\(A\).
\def\ProblemHeadingText{Problem 3.}
\item
Let \(\|\cdot\|_p\) for \(1\leq p\leq\infty\) denote the norm on \(m\times n\) matrices induced by the \(\ell^p\)-norm. For any positive~\(p\) and~\(q\) with \(p^{-1}+q^{-1}=1\), show that
\[
\|A\|_2^2\leq \|A\|_p\,\|A\|_q,\qquad\text{for all }A\in\mathbb C^{m\times n}.
\]
 Hint: Recall that for a Hermitian matrix~\(H\), \(\|H\|_2\) is the absolute largest eigenvalue. As a result, \(\|H\|_2\leq\|H\|\) for any matrix norm induced by a vector norm.
\def\ProblemHeadingText{Problem 4.}
\item
Prove Gershgorin’s theorem: If \(A\in\mathbb C^{n\times n}\), then each eigenvalue of~\(A\) lies in the union of the disks
\[
D_i=\left\{\lambda\in\mathbb C:|\lambda-a_{ii}|\leq\sum_{j\ne i}|a_{ij}|\right\}.
\]
 Moreover, if~\(m\) disks form a connected domain that is disjoint from the other \(n-m\) disks, then there are precisely~\(m\) eigenvalues of~\(A\) within this domain.
\def\ProblemHeadingText{Problem 5.}
\item
Let~\(P\) and~\(Q\) be two \(m\times m\) orthogonal projectors. Prove that \(\|P-Q\|_2\leq 1\).
\def\ProblemHeadingText{Problem 6.}
\item
Assume that \(g(x)\) is continuously differentiable on \([a,b]\), that \(g([a,b])\subseteq[a,b]\), and that
\[
\lambda=\operatorname{Max}_{a\leq x\leq b}|g'(x)|<1.
\]
 Prove that the following are true:
\begin{examsubparts}
\item[\textnormal{(i)}]
\(x=g(x)\) has a unique solution~\(\alpha\) in \([a,b]\).
\item[\textnormal{(ii)}]
For any initial choice \(x_0\) in \([a,b]\), the iteration \(x_{n+1}=g(x_n)\) has the property that \(\lim_{n\to\infty}x_n=\alpha\).
\item[\textnormal{(iii)}]
\[
\lim_{n\to\infty}\frac{\alpha-x_{n+1}}{\alpha-x_n}=g'(\alpha).
\]
\end{examsubparts}
\def\ProblemHeadingText{Problem 7.}
\item
Assume that \(f(x)\), \(f'(x)\), and \(f''(x)\) are continuous in \([a,b]\), and that for some \(\alpha\in(a,b)\), we have \(f(\alpha)=0\) and \(f'(\alpha)\ne 0\). Show that if \(x_0\) is chosen close enough to~\(\alpha\), the iterates
\[
x_{n+1}=x_n-\frac{f(x_n)}{f'(x_n)}
\]
 converge to~\(\alpha\). Moreover,
\[
\lim_{n\to\infty}\frac{\alpha-x_{n+1}}{(\alpha-x_n)^2}=-\frac{f''(\alpha)}{2f'(\alpha)}.
\]
\def\ProblemHeadingText{Problem 8.}
\item
Let \(x_0,x_1,x_2,\ldots,x_n\) be distinct real numbers and let~\(f\) be a given real-valued function with \(n+1\) continuous derivatives. Let~\(I\) be an interval containing all of the \(x_k\), and~\(t\). Show that there is a \(\xi\in I\) such that
\[
f(t)-\sum_{k=0}^n f(x_k)l_k(t)=\frac{(t-x_0)(t-x_1)\cdots(t-x_n)}{(n+1)!}f^{(n+1)}(\xi),
\]
 where \(l_k\) is the Lagrange interpolating polynomial which is 1 at \(x_k\) and 0 at the other \(x_j\), \(j\ne k\).
\def\ProblemHeadingText{Problem 9.}
\item
Let \(I_n(f)=\sum_{j=0}^n w_{j,n}f(x_{j,n})\) be a sequence of approximations to \(I(f)=\int_a^b f(x)\,dx\). Let~\(F\) be a family of functions which is dense in \(C[a,b]\), and suppose that (a) \(I_n(f)\to I(f)\) for all \(f\in F\), and (b) \(B=\sup_n\sum_{j=0}^n|w_{j,n}|<\infty\). Show that \(I_n(f)\to I(f)\) for all \(f\in C[a,b]\).
\def\ProblemHeadingText{Problem 10.}
\item
Consider the differential equation \(y'=f(t,y)\), \(y(0)=y_0\), where \(|f(t,y_1)-f(t,y_2)|\leq K|y_1-y_2|\) for all \(y_1\) and \(y_2\in\mathbb R\). Show that a multistep method
\[
y_{n+1}=\sum_{j=0}^p a_jy_{n-j}+h\sum_{j=-1}^p b_jf(x_{n-j},y_{n-j})
\]
 has local truncation error \(\tau(h)=O(h^m)\) for all \(y(x)\) which are \((m+1)\) times continuously differentiable if and only if
\[
\sum_{j=0}^p a_j=1,
\]
 and
\[
\sum_{j=0}^p(-j)^i a_j+i\sum_{j=-1}^p(-j)^{i-1}b_j=1
\]
 for \(i=1,2,3,\ldots,m\).
\end{examproblems}
\end{document}
